优化并不是先有一个公式，再寻找合适的算法。真正困难的第一步发生得更早：现实中有许多互相冲突的愿望，我们必须说明什么可以改变、什么不能违反，以及究竟怎样比较两个方案。

本单元先完成一次完整的建模：

\begin{enumerate}
\tightlist
\item
  把现实选择翻译成优化问题：把变量、目标、约束、可行域和最优值放回同一个故事中，并区分不可行、无界与最优值不取得。
\item
  局部更好为何还不够：展示一般优化中的局部陷阱，由此提出凸优化真正承诺的事情：在适当的几何结构下，局部证据可以升级为全局结论。
\end{enumerate}

一个求解器只能回答被写下来的问题；如果变量含义、单位、约束方向或目标本身写错，算法越精确，只会越精确地回答错误问题。

\subsection{一个小型的开放决策}

考虑一个小作坊主选择两种产品产量的简化模型。产品 A 单位利润 3 万元，产品 B 单位利润 5 万元。每天总工时 8 小时：A 每件需要 1 小时，B 每件需要 2 小时。原料限制：A 每件用 3 吨，B 每件用 2 吨，总原料 18 吨。令 \(x_A,x_B\) 为日产量，问题为

\[
\max\;3x_A+5x_B,\qquad
x_A+2x_B\le8,\;3x_A+2x_B\le18,\;x_A,x_B\ge0.
\]

这个线性规划的最优解是生产 A 5 件、B 1.5 件，日利润 22.5 万元。但 B 的小数产量提示了模型与现实的距离：假如产品不可分割，需补充整数限制，问题将从线性规划变为整数规划，复杂度剧变。这一个小例包含了建模-求解-验证-修改的完整循环。

\subsection{怎么读这一章}

按真实工作的顺序读。先用第一篇检查问题有没有写对：哪些量能选、哪些规则不能违反、什么才叫更好。再用第二篇检查答案有没有资格被相信：它只是附近走不动，还是已经有全局保证。这里暂时不需要记凸性的所有判别法，先看清建模与认证是两个不同环节。若想回顾上一部分怎样借助额外结构获得更强结论，可回看 \hyperref[ch:069]{``复分析与留数：奇点决定函数的深层结构''}。

\subsection{本章篇目}

以下篇目按概念依赖排列。

\begin{itemize}
\tightlist
\item \hyperref[sec:09-Convex-Optimization/01-Optimization-Foundations/01]{``把现实选择翻译成优化问题''}；
\item \hyperref[sec:09-Convex-Optimization/01-Optimization-Foundations/02]{``局部更好为何还不够''}。
\end{itemize}
